% =========================================================================
% Kapitel 5 – Data Preparation (CRISP-DM Phase 3)
% Stand: 2026-07-19. Zahlen validiert gegen results/eignungspruefung/ und
% results/demo_modellierung/ (Pipeline-Lauf 2026-07-18).
%
% EINFÜGEN: kompletten Inhalt ab \section{Data Preparation} in die Hauptdatei
% übernehmen. Voraussetzung: listings-Preamble (siehe
% kapitel_5_empirische_analyse.tex, Kopfkommentar) ist in der Hauptdatei
% eingebunden.
%
% QUERVERWEISE: Kapitel-/Abschnittsnummern folgen deiner Gliederung
% (4 = Data Understanding, 6 = Modelling, 6.2 = Baselines, 7 = Evaluation,
% 8.3 = Limitationen). Falls du Labels für diese Kapitel hast, die
% "Kapitel~4"/"Abschnitt~8.3"-Stellen durch \ref{...} ersetzen.
%
% KI-VERZEICHNIS: Fußnote in Abschnitt 5.2 (Exposé-Übernahme, Prompt 2)
% beibehalten; Sessions vom 2026-07-18 laut CLAUDE.md §9 ergänzen.
% =========================================================================

\section{Data Preparation}
\label{sec:data-preparation}

Die Data-Preparation-Phase überführt den in Kapitel~4 beschriebenen
integrierten Rohdatenbestand in den finalen Analysedatensatz. Sie umfasst
drei Schritte: die Bereinigung der Quelldaten einschließlich der Behandlung
fehlender Werte (Abschnitt~\ref{sec:bereinigung}), die räumlich-zeitliche
Aggregation auf die Analyseebene (Abschnitt~\ref{sec:aggregation}) sowie das
Feature Engineering mit der Ableitung der endgültigen Prädiktoren
(Abschnitt~\ref{sec:feature-engineering}). Zwei Prinzipien leiten sämtliche
Schritte: Erstens werden alle Transformationen so gestaltet, dass zu keinem
Zeitpunkt Informationen aus der Zukunft in die Merkmale einfließen
(Leakage-Vermeidung);\footnote{Zur Leakage-Problematik bei zeitstrukturierten
Daten vgl. Bergmeir/Benítez (2012), S.~192\,f.} zweitens erhalten alle drei
zu vergleichenden Verfahren denselben aufbereiteten Datensatz, sodass
spätere Leistungsunterschiede ausschließlich auf die Verfahren selbst
zurückführbar sind.

% =========================================================================
\subsection{Bereinigung und Behandlung fehlender Werte}
\label{sec:bereinigung}

Ausgangsbasis sind 720.258 Einsatzdatensätze des SFFD (2003--2026), von
denen bereits beim Abruf nur Zeilen mit vorhandener Stadtteil-Zuordnung und
Ankunftszeit übernommen wurden. Die Bereinigung setzt an drei Punkten an.
Erstens enthalten die Quelldaten 269 mehrfach gemeldete Einsatznummern
(0,04\,\%, davon 50 vollständig identische Zeilen); diese wurden nach der
Einsatznummer dedupliziert, da es sich um Melde-Artefakte der Quelldaten
handelt, die die Einsatzzählungen systematisch überschätzen würden. Es
verbleiben 719.989 Einsätze. Zweitens wurde die Ausrückzeit als Differenz
aus Ankunfts- und Alarmzeitpunkt berechnet und auf den plausiblen Bereich
von 0 bis 60 Minuten begrenzt, wodurch rund 1,7\,\% der Datensätze
entfielen: Negative Werte deuten auf Uhren- bzw.\ Erfassungsfehler, extreme
Werte auf Protokollartefakte. Die Einsatzzählungen beziehen sich damit auf
qualitätsgefilterte Einsätze; diese Einschränkung wird in Abschnitt~8.3
eingeordnet. Drittens wurden in den baulichen Daten Baujahre außerhalb des
Bereichs 1800--2025 als fehlend markiert.

Über die Zeilenbereinigung hinaus wirkt die Leakage-Vermeidung bereits in
dieser Phase als Gestaltungsprinzip der Datenverknüpfung. Die Zuordnung der
jährlich erhobenen ACS-Merkmale zu den Einsätzen erfolgt strikt
prognostisch: Jeder Einsatz erhält den \emph{letzten verfügbaren}
ACS-Jahrgang (Erhebungsjahr $\leq$ Einsatzjahr). Eine Zuordnung des zeitlich
\emph{nächsten} Jahrgangs würde für frühe Einsatzjahre Informationen aus der
Zukunft einschleusen und die spätere Prognosegüte systematisch
überschätzen. Einzige Ausnahme sind die Einsatzjahre 2003--2008, für die
mangels älterer Erhebung auf den ACS-Jahrgang 2009 zurückgegriffen wird;
für den unten begründeten Hauptanalyse-Zeitraum ist dies ohne Belang. Nach
demselben Prinzip werden die Kriminalitätsmerkmale je Stadtteil nur bis zum
Vorjahr des jeweiligen Einsatzes kumuliert. Die baulichen Merkmale liegen
ausschließlich als Querschnitt des Jahres 2020 vor und gehen daher als
zeitkonstante Strukturmerkmale ein; diese Einschränkung wird in
Abschnitt~8.3 diskutiert.

Quellcode~\ref{lst:bereinigung} belegt die drei zentralen Schritte. Die
Deduplikation behält je Einsatznummer das erste Vorkommen
(\texttt{drop\_duplicates}); die Ausrückzeit wird aus der Differenz der
Zeitstempel in Minuten berechnet und der Datenbestand per boolescher
Filterung auf den plausiblen Bereich eingeschränkt. Die Jahrgangszuordnung
wählt für jedes Einsatzjahr per \texttt{max}-Funktion den größten
ACS-Jahrgang, der das Einsatzjahr nicht überschreitet -- die Bedingung
\texttt{a <= int(y)} kodiert damit unmittelbar das Verbot des Blicks in die
Zukunft; der \texttt{default}-Zweig greift ausschließlich für die
Einsatzjahre vor 2009.

\begin{lstlisting}[caption={Deduplikation, Plausibilitätsfilter und strikt
prognostischer ACS-Join (Ausschnitte aus \texttt{pipeline/02\_join.py})},
label=lst:bereinigung]
# Mehrfach gemeldete Einsatznummern aus den Quelldaten entfernen
df = df.drop_duplicates(subset=["incident_number"], keep="first")

# Ausrueckzeit berechnen und auf plausiblen Bereich begrenzen
df["response_time_min"] = (
    df["arrival_dttm"] - df["alarm_dttm"]
).dt.total_seconds() / 60
df = df[(df["response_time_min"] >= 0) & (df["response_time_min"] <= 60)]

# Jeder Einsatz erhaelt den letzten VERFUEGBAREN ACS-Snapshot
# (acs_jahr <= Einsatzjahr) - kein Blick in die Zukunft.
df["acs_year"] = df["year"].apply(
    lambda y: max([a for a in acs_years if a <= int(y)],
                  default=min(acs_years))
)
\end{lstlisting}

Bei den fehlenden Werten sind drei Konstellationen zu unterscheiden.
Erstens lassen sich rund 5,7\,\% der Einsätze keinem ACS-Wert zuordnen, da
für kleine oder erst später ausgewiesene Stadtteile im jeweiligen Jahrgang
kein Census-Tract-Mapping vorliegt: Treasure Island und Mission Bay erhalten
erst ab dem ACS-Jahrgang 2021 Werte, Lakeshore weist keine Baujahrangaben
auf. Die betroffenen Stadtteil-Monate entfallen über den Ausschluss
unvollständiger Zeilen, sodass im Hauptzeitraum 39 der 41 Stadtteile
verbleiben. Zweitens ist die Bildungsvariable maßgeblich: Die ACS-Tabelle
B15003 existiert im Jahrgang 2009 nicht, sodass die Akademikerquote bei
strikt prognostischer Zuordnung erst ab dem Einsatzjahr 2014 zur Verfügung
steht. Anstelle einer Imputation über mehrere Jahre -- insbesondere
anstelle eines Rückwärts-Auffüllens, das erneut Zukunftsinformation
einschleusen würde -- wird die Hauptanalyse auf den Zeitraum ab 2014
begrenzt und der volle Zeitraum ohne dieses Merkmal ergänzend als
Sensitivitätsanalyse berichtet. Drittens wird der Stadtteil McLaren Park als
dokumentiertes Zensus-Artefakt ausgeschlossen: Die ausgewiesene Armutsquote
von 0,90 bei nur 850 Einwohnern einer Parkfläche ist kein valides
sozioökonomisches Merkmal, sondern Folge der Zuschneidung des Zensusbezirks.
Die Robustheit der Ergebnisse gegenüber diesem Ausschluss wird in der
Evaluation geprüft.

% =========================================================================
\subsection{Räumlich-zeitliche Aggregation}
\label{sec:aggregation}

% ---- Übernahme aus Exposé (Fußnote KI-Verzeichnis Prompt 2) ----
Eine zentrale Entwurfsentscheidung betrifft die Granularität: Da die
Feature-Werte auf Stadtteilebene und jährlicher Erhebung vorliegen, würden
sich diese Werte bei einer Auswertung auf Einzeleinsatz-Ebene mehrfach
wiederholen und potenziell zu Pseudo-Signalen führen. Die Daten wurden daher
zeitlich auf Stadtteil~$\times$~Monat aggregiert; Zielgröße für die
Regression ist die Anzahl der Einsätze pro Stadtteil und
Monat.\footnote{In Anlehnung an Anthropic Claude (2026), KI-Verzeichnis
Prompt~2.}
% ---- Ende Übernahme ----

Gegenüber den denkbaren Alternativen ist diese Ebene wie folgt abzuwägen.
Die Einzeleinsatz-Ebene würde die Stichprobe künstlich auf über 700.000
Zeilen aufblähen, obwohl die Prädiktoren nur auf Stadtteilebene variieren;
die Folge wären scheinbar hochsignifikante, tatsächlich aber redundante
Zusammenhänge. Eine Aggregation auf Census-Tract~$\times$~Jahr böte zwar
feinere räumliche Auflösung, scheitert jedoch daran, dass die Einsatzdaten
keine Tract-Zuordnung tragen und eine nachträgliche Geokodierung
fehleranfällig wäre; zugleich würde die Jahresebene die Saisonalität
eliminieren. Die gewählte Ebene Stadtteil~$\times$~Monat ist somit die
kleinste Einheit, auf der Zielgröße und Prädiktoren konsistent vorliegen
und saisonale Dynamik erhalten bleibt.

Technisch wird ein vollständiges Raster aus Stadtteilen und Kalendermonaten
aufgespannt, sodass einsatzfreie Monate als echte Nullen in die Zählung
eingehen; der unvollständige Randmonat wird abgeschnitten. Die je Stadtteil
und ACS-Jahrgang konstanten Merkmale werden in rasterbedingt ergänzten
Monaten ausschließlich \emph{vorwärts} aus der Vergangenheit desselben
Stadtteils aufgefüllt; ein Rückwärts-Füllen wurde aus dem in
Abschnitt~\ref{sec:bereinigung} genannten Grund bewusst unterbunden.
Quellcode~\ref{lst:aggregation} zeigt die Umsetzung: Der
\texttt{groupby}-Schritt zählt die Einsätze je Stadtteil-Monat und übernimmt
die konstanten Stadtteilmerkmale; \texttt{reindex} spannt anschließend das
vollständige Raster auf, wobei neu entstandene Zeilen zunächst leer sind und
ihre Einsatzzahl explizit auf null gesetzt wird (\texttt{fillna(0)}) --
einsatzfreie Monate sind Beobachtungen, keine Datenlücken. Der abschließende
\texttt{ffill}-Aufruf arbeitet je Stadtteil-Gruppe und ausschließlich in
Zeitrichtung.

\begin{lstlisting}[caption={Aggregation auf Stadtteil~$\times$~Monat mit
vollständigem Raster (Ausschnitt aus
\texttt{modellierung/aggregation.py})}, label=lst:aggregation]
agg = (df.groupby(["stadtteil", "jahr", "monat"])
         .agg(anzahl_einsaetze=("einsatz_nummer", "count"),
              **{c: (c, "first") for c in PRAEDIKTOREN})
         .reset_index())

# Vollstaendiges Raster: Monate ohne Einsatz sind echte Nullen
raster = agg.set_index(["stadtteil", "jahr", "monat"]) \
            .reindex(idx).reset_index()
raster["anzahl_einsaetze"] = raster["anzahl_einsaetze"].fillna(0)

# Nur VORWAERTS fuellen - Rueckwaertsfuellen wuerde fehlende Werte
# mit Zukunftsinformation imputieren (Leakage).
raster[PRAEDIKTOREN] = (raster.groupby("stadtteil")[PRAEDIKTOREN]
                              .transform(lambda s: s.ffill()))
\end{lstlisting}

Nach der Aggregation bestätigt sich die in Kapitel~4 beschriebene
Charakteristik der Zielgröße auf der finalen Analyseebene: Die monatliche
Einsatzzahl je Stadtteil ist rechtsschief verteilt (Mittelwert 63,4; Median
45; Maximum 451) und mit einem Dispersionsindex von 61 (Varianz je
Mittelwert) deutlich überdispers, was die Spezifikation der
Negative-Binomial-Regression als Count-Baseline in Abschnitt~6.2
begründet.\footnote{Vgl. Cameron/Trivedi (2013), S.~1\,ff.} Die Zielgröße
der Klassifikation verbleibt demgegenüber bewusst auf der
Einzeleinsatz-Ebene (Brand vs.\ Nicht-Brand; 13,1/86,9\,\%), da die
Einsatzart ein Merkmal des einzelnen Ereignisses ist und die
einsatzbezogenen Zeitmerkmale (Stunde, Wochentag, Nacht) nur dort ihre
Erklärungskraft entfalten.

% =========================================================================
\subsection{Feature Engineering und finaler Analysedatensatz}
\label{sec:feature-engineering}

% ---- Übernahme aus Exposé ----
Im Rahmen des Feature-Engineerings wurden einzelne Prädiktoren aus den
Rohdaten neu berechnet: So wurden etwa der Anteil an Gebäuden mit Baujahr
vor 1940 auf Stadtteilebene sowie ein Index für risikobehaftetes Gewerbe,
auf Basis der Anzahl von Industrie-, Gastronomie- und vergleichbaren
Betrieben, als eigenständige Merkmale konstruiert und dem Modell als
zusätzliche Prädiktoren bereitgestellt.
% ---- Ende Übernahme ----
Die baulichen Merkmale entstammen einem Spatial Join der 155.395 Parzellen
der Land-Use-Datenbank gegen die Stadtteilpolygone (Zuordnung über den
Parzellen-Mittelpunkt, Match-Rate 99,5\,\%). Ergänzend wurden die
sozioökonomischen Quoten (Armuts-, Akademiker- und Leerstandsquote) als
Zähler-Nenner-Verhältnisse im Intervall $[0,1]$ berechnet, die Mediane von
Einkommen und Miete populationsgewichtet von der Tract- auf die
Stadtteilebene aggregiert und die Kriminalitätsanteile als Gewalt- bzw.\
Eigentumsdelikte je Gesamtdelikte gebildet. Die Saisonalität wird über eine
zyklische Kodierung des Monats (Sinus- und Kosinus-Komponente) abgebildet,
die den künstlichen Sprung zwischen Dezember und Januar einer ordinalen
Kodierung vermeidet.

Sämtliche Quoten werden über eine gemeinsame Hilfsfunktion berechnet, die
Quellcode~\ref{lst:features} exemplarisch für die Armutsquote, den
Altbau-Anteil und den Risiko-Gewerbe-Index zeigt. Die Funktion
\texttt{safe\_ratio} erzwingt zunächst numerische Datentypen
(\texttt{errors="coerce"} wandelt nicht interpretierbare Werte in
\texttt{NaN}) und ersetzt Nenner von null oder darunter durch \texttt{NaN},
bevor dividiert wird. Fehlende oder unplausible Grundgesamtheiten führen so
zu einem explizit fehlenden Quotienten statt zu einem Laufzeitfehler oder
einem verzerrten Wert -- die Behandlung dieser fehlenden Werte erfolgt dann
zentral nach den Regeln aus Abschnitt~\ref{sec:bereinigung}.

\begin{lstlisting}[caption={Berechnung der Quoten-Merkmale über eine
robuste Verhältnisfunktion (Ausschnitt aus
\texttt{pipeline/03\_features.py})}, label=lst:features]
def safe_ratio(zaehler, nenner):
    """Quotient in [0, 1]; NaN bei fehlendem/ungueltigem Nenner."""
    z = pd.to_numeric(zaehler, errors="coerce").astype(float)
    n = pd.to_numeric(nenner,  errors="coerce").astype(float) \
          .where(lambda x: x > 0, np.nan)
    return (z / n).round(4)

df["poverty_rate"] = safe_ratio(df["poverty_below"],
                                df["poverty_universe_total"])
df["pct_pre1940"]  = safe_ratio(df["pre1940_count"],
                                df["yrbuilt_count"])
df["pct_high_risk_commercial_area"] = safe_ratio(
    df["high_risk_commercial_area_sqft"], df["total_area_sqft"])
\end{lstlisting}

Über das Exposé hinaus wurden autoregressive Verlaufsmerkmale ergänzt. Den
Anlass gibt der Befund aus Kapitel~4: Bei einer Lag-1-Autokorrelation der
Zielgröße von 0,96 kann kein rein strukturmerkmalbasiertes Modell die naive
Fortschreibung des Vormonatswerts übertreffen. Je Stadtteil werden daher
der Vormonatswert (\texttt{lag\_1}), der Vorjahresmonat (\texttt{lag\_12})
und der gleitende Dreimonatsmittelwert (\texttt{rolling\_mean\_3})
gebildet. Auf ein rohes Jahresmerkmal wird verzichtet, da es bei der
Prognose zukünftiger Monate eine Extrapolation über den Trainingsbereich
hinaus erzwingen würde. Quellcode~\ref{lst:lags} zeigt die Berechnung:
\texttt{shift(1)} bzw.\ \texttt{shift(12)} verschieben die Zeitreihe je
Stadtteil um einen bzw.\ zwölf Monate, sodass jede Zeile ausschließlich
Werte aus ihrer eigenen Vergangenheit erhält. Entscheidend ist die
Reihenfolge beim gleitenden Mittel: Die Reihe wird \emph{erst} um einen
Monat verschoben und \emph{dann} über drei Monate gemittelt -- der aktuelle
Monat fließt somit nie in seine eigenen Merkmale ein, was die
Leakage-Freiheit der Verlaufsmerkmale sicherstellt.

\begin{lstlisting}[caption={Streng vergangenheitsbezogene Lag-Merkmale je
Stadtteil (Ausschnitt aus \texttt{modellierung/demo\_modellierung.py})},
label=lst:lags]
g = daten.groupby("stadtteil")["anzahl_einsaetze"]
daten["lag_1"]          = g.shift(1)          # Vormonat
daten["lag_12"]         = g.shift(12)         # Vorjahresmonat
daten["rolling_mean_3"] = g.transform(        # Mittel der letzten
    lambda s: s.shift(1).rolling(3).mean())   # 3 Monate (ohne t)
\end{lstlisting}

Für die Auswertung werden daraus zwei Merkmalsmengen gebildet: Set~S
umfasst ausschließlich die Struktur- und Saisonmerkmale und dient der
Beantwortung der ersten Unterfrage nach dem Erklärungsbeitrag der
Stadtteilmerkmale; Set~S+L ergänzt die Lag-Merkmale und stellt die faire
Prognoseaufgabe dar, auf der die drei Verfahren in Kapitel~7 verglichen
werden. Tabelle~\ref{tab:praediktoren} führt die finalen Prädiktoren mit
Berechnung und Datenquelle auf.

\begin{table}[htbp]
  \centering
  \caption{Finale Prädiktoren des Analysedatensatzes}
  \label{tab:praediktoren}
  {\small
  \begin{tabular}{llp{5.6cm}l}
    \hline
    \textbf{Merkmal} & \textbf{Gruppe} & \textbf{Berechnung} &
    \textbf{Quelle} \\
    \hline
    Gesamtbevölkerung & Sozioökon. & Wohnbevölkerung des Stadtteils &
      ACS \\
    Medianeinkommen & Sozioökon. & Medianes Haushaltseinkommen,
      populationsgewichtet (USD/Jahr) & ACS \\
    Medianmiete & Sozioökon. & Mediane Bruttomiete,
      populationsgewichtet (USD/Monat) & ACS \\
    Armutsquote & Sozioökon. & Personen unter Armutsgrenze je
      Grundgesamtheit $[0,1]$ & ACS \\
    Akademikerquote & Sozioökon. & Personen mit Bachelor-Abschluss (ab 25
      J.) je Grundgesamtheit $[0,1]$ & ACS \\
    Leerstandsquote & Sozioökon. & Leer stehende je gesamte Wohneinheiten
      $[0,1]$ & ACS \\
    Anteil Gewaltdelikte & Kriminalität & Gewaltdelikte je Gesamtdelikte
      $[0,1]$ & SFPD \\
    Anteil Eigentumsdelikte & Kriminalität & Eigentumsdelikte je
      Gesamtdelikte $[0,1]$ & SFPD \\
    Altbau-Anteil & Baulich & Parzellen mit Baujahr vor 1940 je Parzellen
      mit bekanntem Baujahr $[0,1]$ & Land Use \\
    Wohnnutzungs-Anteil & Baulich & Wohnparzellen je Gesamtparzellen
      $[0,1]$ & Land Use \\
    Risiko-Gewerbe-Index & Baulich & Fläche brandrelevanter
      Gewerbeparzellen (Handel/Gastronomie, Industrie) je Gesamtfläche
      $[0,1]$ & Land Use \\
    Monat (sin, cos) & Saison & Zyklische Kodierung des Kalendermonats
      (2 Merkmale) & abgeleitet \\
    \texttt{lag\_1}, \texttt{lag\_12} & Autoregr. & Einsatzzahl des Vormonats
      bzw.\ Vorjahresmonats & abgeleitet \\
    \texttt{rolling\_mean\_3} & Autoregr. & Gleitendes Mittel der drei
      Vormonate & abgeleitet \\
    \hline
  \end{tabular}}
\end{table}

Modellspezifische Transformationen erfolgen ausschließlich innerhalb der je
Trainingsfenster gefitteten Pipeline, damit keine Kennwerte aus Testdaten in
die Vorverarbeitung einfließen: Für die Ridge Regression werden alle
Prädiktoren z-standardisiert, die Zielgröße wegen der Rechtsschiefe und der
in Kapitel~4 gezeigten Heteroskedastizität als $\log(1+y)$ modelliert und
die Lag-Merkmale ebenfalls $\log(1+x)$-transformiert, da rohe Zählwerte in
einem log-linearen Modell fehlspezifiziert wären; die baumbasierten
Verfahren trainieren unmittelbar auf den unveränderten Zählwerten. Für die
Klassifikation werden das Bataillon sowie die kategorialen Zeitmerkmale
einheitlich für alle drei Verfahren One-Hot-kodiert. Damit gilt
durchgängig: Alle Modelle erhalten identische Zeilen, Merkmale und
Validierungs-Folds; Leistungsunterschiede sind ausschließlich algorithmisch
bedingt.

Den finalen Analysedatensatz fasst Tabelle~\ref{tab:steckbrief} zusammen.
% TODO: Zahlen nach finalem Trainingslauf ggf. aktualisieren
% (Stand: Demo 2026-07-18).

\begin{table}[htbp]
  \centering
  \caption{Steckbrief des finalen Analysedatensatzes}
  \label{tab:steckbrief}
  \begin{tabular}{ll}
    \hline
    \textbf{Merkmal} & \textbf{Ausprägung} \\
    \hline
    Regressionsebene & Stadtteil $\times$ Monat (vollständiges Raster) \\
    Hauptanalyse-Zeitraum & 2014--2026 (2015--2026 nach Lag-Anlauf) \\
    Beobachtungen (Regression) & 5.103 (39 Stadtteile $\times$ 133 Monate) \\
    Zielgröße Regression & Einsätze/Monat (Mittel 63,4; Max.\ 451) \\
    Prädiktoren & 11 Struktur + 2 Saison + 3 Lag \\
    Merkmalsmengen & Set S (Struktur + Saison); Set S+L (zzgl.\ Lags) \\
    Klassifikationsebene & Einzeleinsatz (719.989 Datensätze) \\
    Zielgröße Klassifikation & Brand vs.\ Nicht-Brand (13,1/86,9\,\%) \\
    Ausschlüsse & Duplikate; Ausrückzeit $\notin [0,60]$ min; \\
                & McLaren Park; Stadtteile ohne ACS-/Baujahrdaten \\
    \hline
  \end{tabular}
\end{table}
